\documentclass[10pt,a4paper]{amsart}
\usepackage[margin=19mm]{geometry}
\usepackage{amsmath,amssymb,mathtools}
\usepackage[scr=rsfs,cal=euler]{mathalfa}
\usepackage{xfrac}
\usepackage[unicode,hidelinks,bookmarks=false]{hyperref}
\newcommand{\ie}{i.e.\ }
\newcommand{\cf}{cf.\ }
\newcommand{\resp}{respectively\ }
\newcommand{\Subsection}{\subsection}
\providecommand{\nobreakdash}{\nobreak\hbox{-}}
\setlength{\emergencystretch}{2em}
\multlinegap0pt
\usepackage{url}
\usepackage{xy}
\usepackage{url,color}
\usepackage{algpseudocode,algorithm}
\usepackage{todonotes}
\usepackage{refcheck}
\newcommand{\commS}[2][]{\todo[#1,color=yellow]{S: #2}}
\newcommand{\circc}{\circ^{q}}
\newcommand{\G}{\mathbb G}
\newcommand{\Z}{\mathbb Z}
\newcommand{\g}{\mathcal G}
\newcommand{\m}{\mathbb M}
\newcommand{\rad}{\mathrm{rad}}
\newcommand{\h}{\mathbb H}
\newcommand{\GL}{\mathrm{GL}}
\newcommand{\PGL}{\mathrm{PGL}}
\newcommand{\Tr}{\mathrm{Tr}}
\newcommand{\fqn}{\mathbb{F}_{q^n}}
\newcommand{\F}{\mathbb{F}}
\newcommand{\pp}{\mathcal P}
\newcommand{\fq}{\mathbb{F}_q}
\newcommand{\ord}{\mathrm{ord}}
\newcommand{\ou}{\mathcal O}
\newcommand{\li}{\mathcal L}
\newcommand{\oord}{\mathrm{Ord}}
\newcommand{\order}[1]{\left|#1\right|}
\newtheorem{theorem}{Theorem}[section]
\newtheorem{lemma}[theorem]{Lemma}
\newtheorem{conjecture}[theorem]{Conjecture}
\newtheorem{corollary}[theorem]{Corollary}
\newtheorem{proposition}[theorem]{Proposition}
\newtheorem{definition}[theorem]{Definition}
\newtheorem{example}[theorem]{Example}
\newtheorem{problem}[theorem]{Problem}
\newtheorem{remark}[theorem]{Remark}
\begin{document}
\title[Polynomial identities, central polynomials and cocharacters ]{Polynomial identities, central polynomials and cocharacters of $M_2(F)$ with transpose superinvolution}
\author{Rafael Bezerra dos Santos, Lucas Reis}
\date{}
\maketitle
\noindent\emph{Source and licence.} Polynomial identities, central polynomials and cocharacters of $M_2(F)$ with transpose superinvolution. Version arXiv:2609.20458v1 (CC BY 4.0). This material is distributed under the Creative Commons Attribution 4.0 International licence, http://creativecommons.org/licenses/by/4.0/, as recorded in the arXiv OAI licence metadata for this exact version and shown on the abstract page https://arxiv.org/abs/2609.20458v1. Full rights notice: licenses/arxiv-2609.20458-CC-BY-4.0.txt.\par\smallskip
\noindent\textit{Editorial scope.} Counted unit: the original \texttt{theorem} environment \texttt{thm:main} at source lines 333--345 together with its complete proof at lines 347--371; the delivered document keeps source lines 184--372 so that every referenced label and every citation resolves inside it. Native editable source is the paper's own concatenated TeX; the AMS wrapper is editorial formatting only. No publisher class, upstream script or unrelated artwork is redistributed.
	\begin{remark}\label{consequences}
		\
		
		\begin{itemize}
			\item[1.] $[y_{1,0},z_{2,0}], [y_{1,1},y_{2,1}],[z_{1,1},z_{2,1}],y_{1,0}\circ y_{2,0},y_{1,1}\circ z_{2,1}\in (\mathcal{F}^{(0)})^+$;
			\item[2.] $[y_{1,0},y_{2,0}], [y_{1,1},z_{2,1}], y_{1,0}\circ z_{2,0}, y_{1,1}\circ y_{2,1},z_{1,1}\circ z_{2,1}\in (\mathcal{F}^{(0)})^-$;
			\item[3.] $[y_{1,0},z_{2,1}], [z_{1,0},y_{2,1}],y_{1,0}\circ y_{2,1}, z_{1,0}\circ z_{2,1}\in (\mathcal{F}^{(1)})^+$;
			\item[4.] $[y_{1,0},y_{2,1}], [z_{1,0},z_{2,1}],y_{1,0}\circ z_{2,1}, z_{1,0}\circ y_{2,1}\in (\mathcal{F}^{(1)})^-$.
		\end{itemize}
	\end{remark}
	
	We say that $$f = f(y_{1,0},  \ldots , y_{m,0}, y_{1,1},  \ldots, y_{n,1}, z_{1,0}, \ldots, z_{p,0}, z_{1,1},  \ldots, z_{q,1}) \in \mathcal{F}$$
	is a $*$-identity of a $*$-algebra $A$ if $$f(a_{1,0}, \ldots, a_{m,0},a_{1,1}, \ldots , a_{n,1}, b_{1,0}, \ldots , b_{p,0}, b_{1,1}, \ldots , b_{q,1}) = 0,$$ 
	for all $a_{1,0}, \ldots, a_{m,0} \in (A^{(0)})^{+}$, $a_{1,1}, \ldots, a_{n,1} \in (A^{(1)})^{+}$, $b_{1,0}, \ldots, b_{p,0} \in (A^{(0)})^{-}$ and $b_{1,1}, \ldots , b_{q,1} \in (A^{(1)})^{-}$. In this case, we write $f\equiv 0$ on $A$. The set of all polynomial $*$-identities satisfied by the algebra $A$ 
	$$
	Id^*(A) = \{f \in \mathcal{F}: f \equiv 0 \mbox{ on } A \}
	$$ 
	is a $T^{*}_{2}$-ideal of $\mathcal{F}$, i.e., an ideal of $\mathcal{F}$ invariant under all graded endomorphisms of $\mathcal{F}$ commuting with the superinvolution $*$.
	
	Given a set of $*$-polynomials $S \subseteq \mathcal{F}$, we denote by $\langle S \rangle_{T_2^*}$ the $T_2^*$-ideal of $\mathcal{F}$ generated by the set $S$. Moreover, we say that a set of $*$-polynomials $S'$ is a consequence of $S$ if $S' \subseteq \langle S \rangle_{T_2^*}$. Since char$(F) =0$, each polynomial of $Id^*(A)$ is
	equivalent to a system of multilinear $*$-identities. We denote by
	$$
	P_{n}^{*} = \textrm{span}\{w_{\sigma(1)} \cdots w_{\sigma(n)}: w_{i} \in \{y_{i,0},y_{i,1},z_{i,0},z_{i,1}\}, \sigma \in S_{n}\},
	$$ 
	the space of multilinear $*$-polynomials of degree $n$ in the variables $y_{1,0},  \ldots , y_{n,0}, \\ y_{1,1}, \ldots, y_{n,1}, z_{1,0}, \ldots, z_{n,0}, z_{1,1},  \ldots, z_{n,1}$. Notice that $\dim_FP_n^*=4^nn!$. The dimension of the quotient space
	$$P_n^{*}(A)=\displaystyle \frac{P_n^{*}}{P_n^{*}\cap Id^*(A)}$$ is called the $n$-th $*$-codimension of $A$ and it is denoted by  $c_n^*(A)$.
	
	\begin{remark}
		If $c_n(A)$ denotes the ordinary codimension of a $*$-algebra $A$, then $0\leq c_n(A)\leq c_n^*(A) \leq 4^nc_n(A)$. As a consequence, by \cite{Regev}, if $A$ satisfies a nontrivial ordinary polynomial identity, then $c_n^*(A)$ is exponentially bounded.
	\end{remark}
	
	In \cite{Ninni}, Ioppolo proved the following relations about the sequence of $*$-codimensions of a $*$-algebra. 
	
	\begin{theorem}
		Let $F$ be a field of characteristic zero and $A$ be a $*$-algebra over $F$. Then, there exist constants $C_1>0, C_2, t_1, t_2$ and an integer $d\geq 0$ such that $$C_1n^{t_1}d^n\leq c_n^*(A)\leq C_2n^{t_2}d^n.$$
	\end{theorem}
	
	The integer $d\geq 0$ in the previous theorem is called $*$-exponent of $A$ and it is denoted by exp$^*(A)$.
	
	Now, let $n\geq 1$ and consider $n=n_1+n_2+n_3+n_4$, where each $n_i$ is a nonnegative integer. This is a composition of $n$ and we write $\langle n\rangle=(n_1,n_2,n_3,n_4)$. Given a composition $\langle n\rangle=(n_1,n_2,n_3,n_4)$ of $n$, let $P_{\langle n\rangle}\subseteq P_n^*$ be the vector space of multilinear $*$-po\-ly\-no\-mi\-als in which $n_1$ variables are symmetric of homogeneous degree 0, $n_2$ variables are skew of homogeneous degree 0, $n_3$ variables are symmetric of homogeneous degree 1 and $n_4$ variables are skew of homogeneous degree 1. It is clear that there exist $\binom{n}{{\langle n\rangle}}=\frac{n!}{n_1!n_2!n_3!n_4!}$ subspaces of $P_n^*$ that are isomorphic to $P_{\langle n\rangle}$ and, as a consequence, $$ P_n^*\cong \bigoplus_{{\langle n\rangle}}\binom{n}{{\langle n\rangle}}P_{{\langle n\rangle}},$$  where the sum runs over all compositions $\langle n\rangle$ of $n$. 
	
	Therefore, if $P_{\langle n\rangle}(A)=\dfrac{P_{\langle n\rangle}}{P_{\langle n\rangle}\cap Id^*(A)}$ and $c_{\langle n\rangle}(A)=\dim_F(P_{\langle n\rangle}(A))$, then \begin{equation}\label{codim}
		c_n^*(A)=\displaystyle \sum_{\langle n\rangle}\binom{n}{{\langle n\rangle}}c_{\langle n\rangle}(A).
	\end{equation}
	
	Given a composition $\langle n\rangle=(n_1,n_2,n_3,n_4)$ of $n$, we denote by $S_{\langle n \rangle}=S_{n_1}\times \cdots \times S_{n_4}$ the direct product of four symmetric groups $S_{n_i},1\leq i\leq 4$. Given $\langle n \rangle$, the group $S_{\langle n \rangle}$ acts on the left on $P_{\langle n \rangle}$ as follows: $S_{n_1}$ permutes the variables $y_{1,0}, \ldots, y_{n_1,0}$, $S_{n_2}$ permutes the variables $z_{1,0},\ldots, z_{n_2,0}$, and so on. In this way, $P_{\langle n \rangle}$ becomes a left $S_{\langle n \rangle}$-module. Notice that $P_{\langle n\rangle}\cap Id^*(A)$ is invariant under this action and so $P_{\langle n\rangle}(A)$ is a left $S_{\langle n\rangle}$-module. Its $S_{\langle n\rangle}$-character is denoted by $\chi_{\langle n\rangle}(A)$ and is called the $\langle n\rangle$-cocharacter of $A$.

	Given a composition $\langle n\rangle$ of $n$, a multipartition of $\langle n\rangle$, denoted by $\langle \lambda \rangle\vdash \langle n\rangle$, is a sequence of partitions $\langle \lambda\rangle=(\lambda(1),\ldots, \lambda(4))$ such that $\lambda(i)\vdash n_i,i=1,\ldots, 4$. It is well know that there exists a one-to-one correspondence between partitions of $n$ and irreducible $S_n$-characters. Thus, if $\lambda\vdash n$ , we denote by $\chi_{\lambda}$ the corresponding irreducible $S_n$-character.
	Since $F$ is a field of characteristic zero, by complete reducibility, one can write \begin{equation}\label{ncocharacter}
		\chi_{\langle n\rangle}(A)=\displaystyle \sum_{\langle \lambda \rangle\vdash \langle n\rangle} m_{\langle \lambda \rangle}\chi_{\langle \lambda \rangle},
	\end{equation} where $m_{\langle \lambda \rangle}$ is the corresponding multiplicity of the irreducible $S_{\langle n \rangle}$-character $\chi_{\langle \lambda \rangle}=\chi_{\lambda(1)}\otimes \cdots \otimes \chi_{\lambda(4)}$. In order to compute the multiplicities $m_{\langle \lambda \rangle}$ of the
	irreducible characters appearing in (\ref{ncocharacter}) we use the
	re\-presentation theory of the general linear group $GL_n$.

	Let $F_{n,*}=F\langle y_{1,0},\ldots y_{n,0},z_{1,0},\ldots z_{n,0}, y_{1,1},\ldots y_{n,1}, z_{1,0},\ldots z_{n,1}\rangle$ be the free associative $*$-algebra of rank $n$ and consider $U_1=span_F\{{y_{1,0},\ldots y_{n,0}}\}$, $U_2=span_F\{{z_{1,0},\ldots z_{n,0}}\}$, $U_3=span_F\{{y_{1,1},\ldots y_{n,1}}\}$ and $U_4=span_F\{{z_{1,1},\ldots z_{n,1}}\}.$ There is a natural left action of the group $GL(U_1)\times \cdots\times GL(U_4)\cong GL_n^4$ on the space $U_1\oplus \cdots\oplus U_4$ and we can extend this action diagonally to get an action on $F_{n,*}$. Notice that, if $A$ is a $*$-algebra, then $F_{n,*}\cap Id^*(A)$ is invariant under this action.
	
	Considering $F^m_{n,*}$ the space of all multihomogenous $*$-polynomials of degree $m$ in $F_{n,*}$, the quotient space $F^m_{n,*}(A)=\dfrac{F^m_{n,*}}{F^m_{n,*}\cap Id^*(A)}$ is a $GL_n^4$-module and its $GL_n^4$-character is denoted by $\psi_n^{*}(A)$.
	

	
	In what follows, the reader may consult \cite{livro} for formal definitions and proofs of the results. There exists a one-to-one correspondence between $GL_n^4$-characters and multipartitions $\langle \lambda\rangle\vdash \langle n\rangle$, where each $\lambda(i)$ has at most $m$ parts. Therefore,\begin{equation}\label{cocharacter}
	\psi_n^*(A)=	\displaystyle \displaystyle{\sum_{ \tiny{\begin{array}{ccc}
						\langle \lambda \rangle \vdash \langle n\rangle\\
						h(\lambda(i))\leq m
		\end{array}}}} \bar{m}_{\langle \lambda \rangle}\psi_{\langle \lambda \rangle},
	\end{equation} where $\bar{m}_{\langle \lambda \rangle}$ is the corresponding multiplicity of the irreducible $GL_n^4$-character $\psi_{\langle \lambda \rangle}$ associated to the multipartition $\langle \lambda \rangle$ and $h(\lambda(i))$ denotes the height of the Young diagram corresponding to $\lambda(i)$ for each $1\le i\le 4$.
	
	\ A Young multitableau $T_{\langle \lambda \rangle}$ of shape $\langle \lambda \rangle$ is a 4-tuple $(T_{\lambda(1)},\ldots, T_{\lambda(4)})$ of Young tableaux. Each multitableau $T_{\langle \lambda \rangle}$ can be associated to a polynomial $f_{
	T_{\langle \lambda \rangle}},$ called {\em highest weight vector}. Moreover, an irreducible submodule of $F^m_{n,*}(A)$ corresponding to the multipartition $\langle \lambda \rangle$ is generated by a non-zero polynomial $f_{\langle \lambda \rangle}$ which can be written as a linear combination of highest weight vectors $f_{
	T_{\langle \lambda \rangle}}.$
	
	\begin{remark}\label{multiplicities}
		
		\
		
		\begin{itemize}
			\item[1.] We have that  $m_{\langle \lambda \rangle}=\bar{m}_{\langle \lambda \rangle}$, for all multipartitions $\langle \lambda \rangle$ such that $h(\lambda(i))\leq m, i=1,\ldots,4$;
			
			\item[2.] For the decomposition in Eq.~\eqref{cocharacter}, we have $\bar{m}_{\langle \lambda \rangle}\neq 0$ if and only if there exists a multitableau $T_{\langle \lambda \rangle}$ such that the corresponding highest weight vector satisfies $f_{T_{\langle \lambda \rangle}}\notin Id^{*}(A)$. Moreover, $\bar{m}_{\langle \lambda \rangle}$ is the maximal
			number of linearly independent highest weight vectors $f_{T_{\langle \lambda \rangle}}$ in $F^m_{n,*}(A)$;
			\item[3.] If $A$ is a $*$-algebra, $A=(A^{(0)})^+ +(A^{(0)})^- + (A^{(1)})^+ + (A^{(1)})^-,$ and $\dim_F((A^{(i)})^{\epsilon})=d_i^{\epsilon}, i\in \{0,1\},\epsilon\in\{+,-\}$, then $\bar{m}_{\langle \lambda \rangle}= 0$ if $h(\lambda(1))>d_0^+, h(\lambda(2))>d_0^-,h(\lambda(3))>d_1^+$ and $h(\lambda(4))>d_1^-.$
		\end{itemize}
		
	\end{remark}
	
	Next, we will deal with central polynomials. Given an algebra with superinvolution $A$, a $*$-polynomial $f\in \mathcal{F}$ is a central $*$-polynomial for $A$ if it has no constant term and any admissible evaluation on $f$ at elements of $A$ belongs to the center of $A$. It is clear that a $*$-polynomial identity of $A$ is a central $\ast$-polynomial of $A$. The set $$Id^{*,z}(A)=\{f\in \mathcal{F}:\mbox{ $f$ is a central $*$-polynomial for $A$}\}$$ is not, in general, a $T_2^*$-ideal of $\mathcal{F}$. However, it is a $T_2^*$-subspace of $\mathcal{F}$, i.e. a subspace invariant under all graded endomorphisms of $\mathcal{F}$ commuting with the superinvolution $*$. If $S\subset \mathcal{F}$, we denote by $\langle S\rangle^{T_2^*}$ the $T_2^*$-subspace generated by $S$. Notice that $\langle S\rangle^{T_2^*}$ is generated, as a subspace, by $f(g_1,\ldots,g_n)$, where $f\in S$ and each $g_i$ has the same homogeneous degree and symmetry
as the variable it replaces.

     \begin{remark}\label{rem:center}
         For $n\geq 1$, the center of $M_n(F)$ is isomorphic to $F$. More precisely, $Z(M_n(F))=\{\alpha I_n:\alpha \in F\},$ where $I_n$ denotes the $n\times n$ identity matrix.
     \end{remark}
     
	 
	 In the next section, our main goal is to determine the sets $Id^*(A)$ and $Id^{*,z}(A)$, the explicit decomposition of $\chi_{\langle n\rangle}(A)$ and to provide asymptotics for $c_n^{*}(A)$, when $A$ is the algebra $M_{1,1}(F)$ endowed with transpose superinvolution.
	 
	We end this section with some relevant auxiliary lemmas. In the next results, $F$ is a field of arbitrary characteristic.
	
		\begin{lemma}\label{lem:vanish}
		Let $F$ be an infinite field and let $f$ be a polynomial over $F$ in $n$ commuting variables. If $f$ vanishes entirely on $F^n$, then $f$ equals the zero polynomial.
	\end{lemma}
	
	Given $\mathcal{F}=F \langle X|* \rangle$, we denote by $\mathcal{F}^C=F[X|*]$ the free associative and commutative $*$-algebra on $X$ over $F$.
	
	\begin{definition}
		
		Let $f\in \mathcal{F}$. We denote by $f^C\in \mathcal{F}^C$ the polynomial induced by $f$ through the relations $x_ix_j=x_jx_i$ for each $x_k\in\{y_{k,0},z_{k,0},y_{k,1},z_{k,1}\}$ with  $k=i,j$.
	\end{definition}
	For instance, if $f(y_{1,0}, z_{1,0})=y_{1,0}z_{1,0}+z_{1,0}y_{1,0}\in \mathcal{F}$, then $f^C(y_{1,0}, z_{1,0})=2y_{1,0}z_{1,0}\in \mathcal{F}^C$. The polynomial $f^C$ can be zero even if $f\in \mathcal F$ is nonzero: $f(y_{1, 0}, z_{1, 0}, z_{1, 1})=y_{1,0}z_{1,0}z_{1, 1}-z_{1,0}z_{1, 1}y_{1,0}\in \mathcal F$ satisfies $f(y_{1, 0}, z_{1, 0}, z_{1, 1})^C=0\in \mathcal F^C$. The following lemma, crucial in the proofs of our main results, provides a sufficient condition under which the vanishing of $f^C$ implies the vanishing of $f$.
	
\begin{lemma}\label{lem:com}
		Let $F$ be a field and let $f\in \mathcal{F}$ be a polynomial not admitting two monomials $M_1, M_2$ with $(M_1)^C=(M_2)^C$. If $f^C\in \mathcal{F}^C$ is the zero polynomial, the same holds for 
		$f$. In particular, if $F$ is infinite, $f$ contains $n$ distinct variables and $f^C$ vanishes entirely on $F^n$, then $f$ is the zero polynomial.
	\end{lemma}
	\begin{proof}
		The first statement follows directly by induction on the number of monomials in the expansion of $f$. The second statement follows directly by the first one and Lemma~\ref{lem:vanish}.
	\end{proof}
	
	
	
	
	
	
	
	
	
	\section{$*$-polynomial identities, $*$-codimension and the $\langle n \rangle$-cocharacter}
	For the rest of the paper, $\mathbb M$ will denote the $*$-algebra $M_{1,1}(F)$ endowed with transpose superinvolution. In this section, we will present the generators of the $T_2^*$-ideal $Id^*(\mathbb M)$, compute the $*$-codimension $c_n^*(\mathbb M)$ and exhibit the decomposition of the $\langle n\rangle$-cocharacter of $\mathbb M$. In the proof of our main results we will work with special monomials in $\mathcal F$. To keep the notation concise and avoid long expressions, we introduce the following notation.
    
	\begin{definition}\label{def:basic}

    \
    
		\begin{enumerate}
		    \item For subsets $S, T\subseteq \{1, \ldots, n\}$, we set $$(YZ)_{S, T}=y_{s_1, 0}\cdots y_{s_j, 0}z_{t_1, 0}\cdots z_{t_k, 0},$$ 
		where $(s_i)_{1\le i\le j}$ and $(t_i)_{1\le i\le k}$ denote the increasing sequence of the elements in $S$ and $T$, respectively. 

        \item For $U, V\subseteq \{1, \ldots, n\}$ with $|U|=|V|=j>0$, set 
		$$(Y|Z)_{U, V}=y_{u_1, 1}z_{v_1, 1}\cdots y_{u_j, 1}z_{v_j, 1},$$
		where $(u_i)_{1\le i\le j}$ and $(v_i)_{1\le i\le j}$ denote the increasing sequence of the elements in $U$ and $V$, respectively. We define $(Z|Y)_{U, V}$ in a similar way.

        \item For $U, V\subseteq \{1, \ldots, n\}$ with $|U|=j+1$ and $|V|=j\ge 0$, set 
		$$\{Y|Z\}_{U, V}=y_{u_1, 1}z_{v_1, 1}\cdots y_{u_j, 1}z_{v_j, 1}y_{u_{j+1}},$$
		where $(u_i)_{1\le i\le j+1}$ and $(v_i)_{1\le i\le j}$ denote the increasing sequence of the elements in $U$ and $V$, respectively. We define $\{Z|Y\}_{U, V}$ in a similar way when $|V|=j+1$ and $|U|=j\ge 0$.
		\end{enumerate}
	\end{definition}
	
	We recall that $(\mathbb M^{(0)})^+=span_F\{e_{11}+e_{22}\}$, $(\mathbb M^{(0)})^-=span_F\{e_{11}-e_{22}\}$, $(\mathbb M^{(1)})^+=span_F\{e_{21}\}$ and $(\mathbb M^{(1)})^-=span_F\{e_{12}\}$. In what follows, we provide a set of generators for $Id^*(\mathbb M)$. 
	
	

	\begin{theorem}\label{thm:main}
		Let $F$ be a field of characteristic zero. Then the $T_2^*$-ideal $I=Id^*(\mathbb M)$ is generated by the following polynomials:
		\begin{enumerate}
			\item $[y_{1, 0}, x]$ with $x\in \{y_{2,0},
			y_{1, 1}, z_{1, 0}, z_{1, 1}\}$;
			\item $y_{1, 1}y_{2, 1}$;
			\item $z_{1, 1}z_{2, 1}$;
			\item $[z_{1, 0}, z_{2, 0}]$;
			\item $z_{1, 0}\circ z_{1, 1}$;
			\item $z_{1, 0}\circ y_{1, 1}$.
			% \item $y_{1, 1}z_{1, 1}y_{2, 1}=y_{2, 1}z_{1, 1}y_{1, 1}$.
		\end{enumerate}
	\end{theorem}

	\begin{proof}
		Let $J$ be the $T_2^*$-ideal generated by the above polynomials. It is direct to verify that $J\subseteq I$. Now let $f\in I$. Since $F$ has characteristic $0$, we can assume that $f$ is multilinear of degree $n\ge 1$. From the identities in (1), (4), (5) and (6), we see that $f\equiv P\pmod J$, where $P$ is a multilinear polynomial of degree $n$ with the following property: 
		each monomial of $P$ is of the form $\alpha X_0X_1$, where
		$X_1$ is a monomial of degree $k\ge 0$ in the variables of homogeneous degree one and
		$X_0=(YZ)_{S, T}$ for some $S, T\subseteq \{1, \ldots, n\}$ with $|S|+|T|=n-k$. Moreover, from the identities (2) and (3), we can also assume that the monomial $X_1$ alternates between the two types of variables of homogeneous degree one. Since 
		$$y_{1, 1}\circ [z_{1, 1}, y_{2, 1}]=y_{1, 1}z_{1, 1}y_{2, 1}-y_{1, 1}y_{2, 1}z_{1, 1}+z_{1, 1}y_{2, 1}y_{1, 1}-y_{2, 1}z_{1, 1}y_{1, 1},$$
		identities (2) and (3) imply 
		$$y_{1, 1}\circ [z_{1, 1}, y_{2, 1}]\equiv y_{1, 1}z_{1, 1}y_{2, 1}-y_{2, 1}z_{1, 1}y_{1, 1}\pmod J.$$
		Now, by Remark \ref{consequences}, we have that $ [z_{1, 1}, y_{2, 1}]\in (\mathcal{F}^{(0)})^-$. Hence identity (6) entails that $y_{1, 1}\circ [z_{1, 1}, y_{2, 1}]\in J$. In conclusion,  
		$$y_{1, 1}z_{1, 1}y_{2, 1}\equiv y_{2, 1}z_{1, 1}y_{1, 1}\pmod J.$$
		In a similar way we obtain $z_{1, 1}y_{1, 1}z_{2, 1}\equiv z_{2, 1}y_{1, 1}z_{1, 1}\pmod J.$
		From these two consequences we can also assume that, among the variables of the same type in $X_1$, their indices are ordered in increasing order. 
		
		With Definition~\ref{def:basic} in mind, these observations lead to the following: the polynomial $P$ is of the form $P_0+P_1+P_2+P_3+P_4$, where $P_0$ only contains monomials of the form $\alpha_0\cdot (YZ)_{S, T}$, $P_1$ only contains monomials of the form $\alpha_1\cdot (YZ)_{S, T}\cdot (Y|Z)_{U, V}$, $P_2$ only contains monomials of the form $\alpha_2\cdot (YZ)_{S, T}\cdot (Z|Y)_{U, V}$, $P_3$ only contains monomials of the form $\alpha_3\cdot (YZ)_{S, T}\cdot \{Y|Z\}_{U, V}$ and $P_4$ only contains monomials of the form $\alpha_4\cdot (YZ)_{S, T}\cdot  \{Z|Y\}_{V, U},\alpha_i\in F$. We further write $P_i=P_{i, +}+P_{i, -}$, where $P_{i, +}$ and $P_{i, -}$ are the sums of the monomials of $P_i$ containing an even and odd number of variables $z_{i, 0}$, respectively. We want to prove that each $P_{i, \pm}$ vanishes. Our idea is to take suitable evaluations at the variables and employ Lemma~\ref{lem:com}. From construction, it is clear that the polynomials $P_{i, \pm}$ and $P_{i, +}\pm P_{i, -}$ do not contain two monomials $M_1, M_2$ with $M_1^C=M_2^C$. 
		
		From hypothesis, $P$ is an identity. Take $y_{i, 1}= z_{i, 1}=0$, $y_{i, 0}=a_i(e_{11}+e_{22})$ and $z_{i, 0}=b_i (e_{11}-e_{22})$ with $a_i, b_i\in F$. We obtain
		$$(Q_{0, +}+Q_{0, -})e_{1 1}+(Q_{0, +}-Q_{0, -})e_{2 2}=0,$$
		where $Q_{0, \pm}\in F$ is the image of $P_{0, \pm}^C$ at the elements $a_i, b_i$. Since $e_{11}$ and $e_{22}$ are linearly independent, we obtain $Q_{0, +}+Q_{0, -}=Q_{0, +}-Q_{0, -}=0$ and so $Q_{0, +}=Q_{0, -}=0$. From Lemma~\ref{lem:com}, both $P_{0, +}$ and $P_{0, -}$ are zero. In conclusion, $P_0$ equals zero.
		
		Now take $y_{i, 0}=a_i(e_{11}+e_{22}), z_{i, 0}=b_i(e_{11}-e_{22}), y_{i, 1}=c_i e_{21}$ and $z_{i, 1}=d_i e_{12}$ for each $1\le i\le n$.
		Similarly to the previous computation, we obtain 
		$$(Q_{1, +}-Q_{1, -})e_{22}+(Q_{2, +}+Q_{2, -})e_{11}+(Q_{3, +}-Q_{3, -})e_{21}+(Q_{4, +}+Q_{4, -}) e_{12}=0,$$
		where $Q_{k, \pm}\in F$ is the image of $P_{k, \pm}^C$ at the elements $a_i, b_i, c_i, d_i$. 
		Again, as the elements $e_{ij}$ are linearly independent, we obtain $Q_{k, +}+(-1)^kQ_{k, -}=0$ for each $1\le k\le 4$. From Lemma~\ref{lem:com}, we obtain $P_{k, +}+(-1)^kP_{k, -}=0$. From construction, the monomials in $P_{k, +}$ cannot cancel with any monomial in  $P_{k, -}$ so the last equality implies that $P_{k, +}=P_{k, -}=0$. In conclusion, each $P_k$ equals zero and so $P$ equals zero. Since $f\equiv P\pmod J$, we obtain $f\in J$ and then $I\subseteq J$, concluding the proof.
	\end{proof}

\begin{thebibliography}{99}
\bibitem[Ninni]{Ninni} A. Ioppolo, The exponent for superalgebras with superinvolution, Linear Algebra Appl. {\bf 555} (2018), 1--20; MR3834188
\bibitem[Regev]{Regev} A. Regev, Existence of identities in $A\otimes B$, Israel J. Math. {\bf 11} (1972), 131--152; MR0314893
\bibitem[livro]{livro} V.~S. Drensky, {\it Free algebras and PI-algebras}, Springer, Singapore, 2000; MR1712064
\end{thebibliography}
\end{document}
